\section{Conclusion}

Fumo Lab's vision is a compound intelligence system whose progress is not
limited to the capability curve of a single model. Fusion advances through two
coupled forces: the model ecosystem supplies stronger and more diverse
intelligence, while the allocation layer learns how to turn that supply into
better outcomes.

The need for this system is becoming more pronounced. Agentic work turns a
request into a trajectory in which the capability and evidence required next
are revealed through execution. At the same time, the model ecosystem is
becoming increasingly heterogeneous. Request-level routing can choose an
answering model, but it cannot fully determine in advance which uncertainty
will become load-bearing as the work unfolds.

Fusion addresses this gap by organizing computation around decision-relevant
uncertainty. It treats models and tools as conditional sources of evidence,
preserves one canonical state, and reserves answers and actions for one
authoritative commitment path. This allows heterogeneous capabilities to
function as one coherent intelligence across both one-shot problems and
long-running workflows.

The immediate result is the possibility of a Pareto improvement: frontier-level
performance with materially lower frontier-compute consumption. The longer-term
result is a system that learns from its own operation. Allocation decisions,
producer behavior, synthesizer use, and downstream outcomes improve the task
representation, capability map, and allocation policy that govern future work.

This is Fusion's bounded form of system-level recursive improvement. Model
progress creates more intelligence; allocation progress determines how much of
it becomes useful; experience improves the allocation of the next task. Fumo
Lab's ambition is to make those forces compound.
